\begin{center}
  {\fontsize{14pt}{16.8pt}\bfseries [论文题目]}\\[0.65em]
  {\fontsize{14pt}{16.8pt}\bfseries 摘\quad 要}
\end{center}

% 摘要写作指南（完成后删除本注释块）：
%
% 结构：背景与核心困难 → 关键建模定义或结构 → 主方法 →
%        最重要 2-3 个结果（含数字） → 结果规律与决策含义 → 可信边界
%
% 禁止：按问题编号（问题一、问题二……）逐句点名。
%        除非各问方法完全不同，否则不要在摘要里提"问题一/二/三"。
%        禁止用"建立了完整的数学模型体系"等空洞总结句开头。
%
% 示例结构（替换为实际内容）：
% 针对 [核心困难]，本文通过 [关键建模定义] 将问题转化为 [可解形式]。
% 采用 [主方法]，得到 [最重要结果，含数字]。
% 分析表明 [规律与决策含义]。
% 上述结论在 [约束条件] 下成立；若 [关键假设改变]，结论将 [如何变化]。
%
% 示例（应替换为实际数字）：
[背景：核心困难是什么，使这道题非显然]。本文将 [问题的核心困难] 定义为 [建模结构]，
以 [主算法/框架] 求解。得到 [最重要数值结果，如：最优方案使总收益达到 X，较基线提升 Y\%]。
分析发现 [核心规律，如：第 2、3 枚烟幕存在边际收益递减，主要瓶颈来自共享航向约束]，
据此 [决策含义]。模型在 [参数范围] 内鲁棒；当 [关键假设] 改变时结论可能随之调整。

\vspace{0.6em}

\noindent{\heiti\bfseries 关键词：}[关键词 1]；[关键词 2]；[关键词 3]；[关键词 4]

\newpage
